% GENERATED FILE. Edit canonical lessons, then run npm run generate:books.
% canonical-source-hash: fnv1a64:7846233774b9ef13
% canonical-lessons: MR-C126-ravivar, MR-C126-tarikh, MR-C126-sadhya, MR-C126-roj, MR-C126-unhala

\chapter{Sunday, Date, At present, Every day, Summer}
\label{ch:mr-sunday-date-at-present-every-day-summer}
\hlchaptermodality{126}

\begin{chapteropening}
\textbf{By the end of this chapter:} \emph{I can say Sunday, a date, at present, every day and summer.}

Complete the last lesson of chapter 126: i can say Sunday, a date, at present, every day and summer.
\end{chapteropening}

\section[ravivār]{\mr{रविवार} (ravivār) — Sunday}

\label{lesson:MR-C126-ravivar}



\begin{quote}
Before the new one: say the Marathi for Friday, then the Marathi for Saturday.
\end{quote}

\subsection*{You'll want to know: \mr{रविवार}}
\textbf{\mr{रविवार}} — \emph{ravivār} — \textquotedblleft{}Sunday\textquotedblright{}.

\begin{grammarlens}[title={What you've built}]
One more word for this chapter.
\end{grammarlens}

\subsection*{Your turn}
\begin{itemize}
\raggedright
  \item \emph{Say it:} \emph{ravivār}
  \item \emph{Say it:} \emph{ravivār}, once more
  \item \emph{Say it:} say \emph{ravivār}
  \item \emph{Recall:} say the Marathi for Friday, then the Marathi for Saturday, then say \emph{ravivār} again
\end{itemize}

\begin{tcolorbox}[breakable,colback=teal!4,colframe=teal!35!black,title={Before you move on}]
What does \mr{रविवार} mean? (\textquotedblleft{}Sunday\textquotedblright{}.) Say it once more.
\end{tcolorbox}
\section[tārīkh]{\mr{तारीख} (tārīkh) — a date}

\label{lesson:MR-C126-tarikh}



\begin{quote}
Before the new one: say the Marathi for Saturday, then the Marathi for Sunday.

Put an event at a time: say \textbf{\mr{पाच} \mr{वाजता}}, at five o'clock.
\end{quote}

\subsection*{You'll want to know: \mr{तारीख}}
\textbf{\mr{तारीख}} — \emph{tārīkh} — \textquotedblleft{}a date\textquotedblright{}.

\begin{grammarlens}[title={What you've built}]
One more word for this chapter.
\end{grammarlens}

\subsection*{Your turn}
\begin{itemize}
\raggedright
  \item \emph{Say it:} \emph{tārīkh}
  \item \emph{Say it:} \emph{tārīkh}, once more
  \item \emph{Say it:} say \emph{tārīkh}
  \item \emph{Recall:} say the Marathi for Saturday, then the Marathi for Sunday, then say \emph{tārīkh} again
\end{itemize}

\begin{tcolorbox}[breakable,colback=teal!4,colframe=teal!35!black,title={Before you move on}]
What does \mr{तारीख} mean? (\textquotedblleft{}A date\textquotedblright{}.) Say it once more.
\end{tcolorbox}
\section[sadhyā]{\mr{सध्या} (sadhyā) — at present, these days}

\label{lesson:MR-C126-sadhya}



\begin{quote}
Before the new one: say the Marathi for Sunday, then the Marathi for a date.
\end{quote}

\subsection*{You'll want to know: \mr{सध्या}}
\textbf{\mr{सध्या}} — \emph{sadhyā} — \textquotedblleft{}at present, these days\textquotedblright{}.

\begin{grammarlens}[title={What you've built}]
One more word for this chapter.
\end{grammarlens}

\subsection*{Your turn}
\begin{itemize}
\raggedright
  \item \emph{Say it:} \emph{sadhyā}
  \item \emph{Say it:} \emph{sadhyā}, once more
  \item \emph{Say it:} say \emph{sadhyā}
  \item \emph{Recall:} say the Marathi for Sunday, then the Marathi for a date, then say \emph{sadhyā} again
\end{itemize}

\begin{tcolorbox}[breakable,colback=teal!4,colframe=teal!35!black,title={Before you move on}]
What does \mr{सध्या} mean? (\textquotedblleft{}At present, these days\textquotedblright{}.) Say it once more.
\end{tcolorbox}
\section[roj]{\mr{रोज} (roj) — every day}

\label{lesson:MR-C126-roj}



\begin{quote}
Before the new one: say the Marathi for a date, then the Marathi for at present.
\end{quote}

\subsection*{You'll want to know: \mr{रोज}}
\textbf{\mr{रोज}} — \emph{roj} — \textquotedblleft{}every day\textquotedblright{}.

\begin{grammarlens}[title={What you've built}]
One more word for this chapter.
\end{grammarlens}

\subsection*{Your turn}
\begin{itemize}
\raggedright
  \item \emph{Say it:} \emph{roj}
  \item \emph{Say it:} \emph{roj}, once more
  \item \emph{Say it:} say \emph{roj}
  \item \emph{Recall:} say the Marathi for a date, then the Marathi for at present, then say \emph{roj} again
\end{itemize}

\begin{tcolorbox}[breakable,colback=teal!4,colframe=teal!35!black,title={Before you move on}]
What does \mr{रोज} mean? (\textquotedblleft{}Every day\textquotedblright{}.) Say it once more.
\end{tcolorbox}
\section[unhāḷā]{\mr{उन्हाळा} (unhāḷā) — summer}

\label{lesson:MR-C126-unhala}



\begin{quote}
Before the new one: say the Marathi for at present, then the Marathi for every day.
\end{quote}

\subsection*{You'll want to know: \mr{उन्हाळा}}
\textbf{\mr{उन्हाळा}} — \emph{unhāḷā} — \textquotedblleft{}summer\textquotedblright{}.

\begin{grammarlens}[title={What you've built}]
One more word for this chapter.
\end{grammarlens}

\subsection*{Your turn}
\begin{itemize}
\raggedright
  \item \emph{Say it:} \emph{unhāḷā}
  \item \emph{Say it:} \emph{unhāḷā}, once more
  \item \emph{Say it:} say \emph{unhāḷā}
  \item \emph{Recall:} say the Marathi for at present, then the Marathi for every day, then say \emph{unhāḷā} again
\end{itemize}

\begin{tcolorbox}[breakable,colback=teal!4,colframe=teal!35!black,title={Before you move on}]
What does \mr{उन्हाळा} mean? (\textquotedblleft{}Summer\textquotedblright{}.) Say it once more.
\end{tcolorbox}
